\section{官僚整顿与继承的脆弱平衡}

1307年海山即位，1311年爱育黎拔力八达即位，1313年元廷恢复科举。这些日期可由《元史》本纪建立年序；科举诏令和制度条格可以说明国家重开了一条取士通道，却不能据此推定各地学校、士人和官场同样受益。以汉文官僚制度解释整个元朝同样不足：多语文书、站赤、军户、寺院和各地区的行政实践并不在同一套材料中显现。

1320年英宗即位，试图调整宫廷与官僚权力；1323年他在南坡遇刺，泰定帝即位；1328年两都之战爆发。权力转换的核心事实可靠，但刺杀的直接策划、每个集团的动机和军队规模，多受胜者记录与后世归罪影响。它们更稳妥地显示，皇位继承与大都、上都、诸王及官僚网络的资源分配始终相互牵连。

1333年妥懽帖睦尔即位，1340年脱脱掌政后推动修史、钞法和行政整顿。修史、货币和机构改革留下较多朝廷文本，但颁行、复申和地方执行必须分开。特别是钞法，发行数额不能当作流通量，也不能直接解释物价和家庭负担；要理解其效果，仍需钱币、契约、地方账册和市场材料。

\section{河工、动员与多中心崩解}

1344年黄河决溢及灾害加深北方的社会和财政压力。灾异记录有明显的政治修辞，不能把一场决溢写成王朝覆亡的单一原因；不过河工、赈济、漕运和税粮确会在同一时期争夺人力与财力。1348年方国珍在浙东海域起事，已显示沿海的地方武装、贸易与中央控制并不能由北方财政单线统摄。

1351年红巾军起事，元廷又发动大规模河工。两件事相连，却不能把被征发者、起事者和受灾者假定为同一群体；不同河段、城市和乡村所受压力须分别考察。1354年脱脱在镇压过程中失势，中央失去了一个协调河工、军队与政治联盟的节点，但地方军阀和起事集团的力量并未因此遵循同一方向。

1356年朱元璋攻取集庆并改应天，1363年鄱阳湖之战改变江南群雄力量对比，1367年明军北伐，1368年进入大都、顺帝北撤。《元史》以明初新王朝的终结视角编纂，关于战役、降附和“失国”的语言必须同地方、碑刻与其他区域材料互校。大都失守标志元朝在中国的中枢退出，却不等于蒙古政治共同体消失；北撤后的政权、仍在运作的地方关系以及被战争改变的人群，需要在“1368年终结”之外继续追问。
